%----------------------------------------------------------------------------------------
% VERGELIJKINGEN EN FORMULES. Alle opgaven zijn eigen werk.
%
% Een hoofdstuk kernvaardigheden (ontwerpregels, deel Kernvaardigheden): beknopte theorie, daarna drill.
% Labels van dit hoofdstuk beginnen met h5-.
%----------------------------------------------------------------------------------------

\mychapter{5}{Vergelijkingen}

\lkr{Herhaling van Eerstegraadsvergelijkingen (3de jaar) en Tweedegraadsfuncties (4de jaar: tweedegraadsvergelijkingen). Nieuw voor sommige leerlingen: vergelijkingen met breuken (bundel 8) en formules omvormen (bundel 9), nodig in de wetenschapsvakken. Bundel 11 veronderstelt Veeltermen. Veronderstelt hoofdstuk 3 (ontbinden) en 4 (breuken).}

\section{Wat je moet kennen}

\begin{theorie}
\begin{minipage}[t]{.47\linewidth}
\vspace{0pt}
\textbf{\textcolor{accent}{Product-nulregel}}

\vspace{1mm}
$\formule{A \cdot B = 0 \iff A = 0 \text{ of } B = 0}$

\vspace{2mm}
\textbf{\textcolor{accent}{Ongelijkheden}}: vermenigvuldig of deel je beide leden door een negatief getal, dan keert het teken om.
\end{minipage}\hfill
\begin{minipage}[t]{.5\linewidth}
\vspace{0pt}
\textbf{\textcolor{accent}{Tweedegraadsvergelijking}} $ax^2 + bx + c = 0$

\vspace{1mm}
$\formule{D = b^2 - 4ac}$ \qquad $\formule{x_{1,2} = \dfrac{-b \pm \sqrt D}{2a}}$

\vspace{1mm}
$D > 0$: twee oplossingen, $D = 0$: één, $D < 0$: geen.
\end{minipage}
\end{theorie}

\begin{werkwijze}
\textbf{Eerste graad:} haakjes weg, noemers weg (maal het kgv), termen met $x$ naar één kant. \textbf{Tweede graad:} herleid op nul. Ontbreekt $c$, zonder dan $x$ af; ontbreekt $b$, isoleer dan $x^2$; anders ontbinden of de discriminant. \textbf{Met breuken:} eerst de bestaansvoorwaarden, en controleer de oplossingen.
\end{werkwijze}

\textbf{Klassieke fouten.} Delen door $x$: uit $x^2 = 3x$ volgt niet alleen $x = 3$, ook $x = 0$. Uit $x^2 = 25$ volgt $x = 5$ \emph{of} $x = -5$. De product-nulregel werkt alleen met $0$ als rechterlid: bij $(x - 2)(x + 1) = 4$ moet je eerst uitwerken.

\voorbeeld{}{Los op: \quad a. $2x^2 = 6x$ \qquad b. $\dfrac{x^2}{x - 3} = \dfrac{9}{x - 3}$\label{vb:h5-vgl}}{%
\begin{stappen}
2x(x - 3) = 0, \text{ dus } x = 0 \text{ of } x = 3 & a. Herleid op nul en zonder $2x$ af. \\
x \neq 3, \quad x^2 = 9, \text{ dus } x = -3 & b. Maal $x - 3$; de oplossing $3$ valt weg door de voorwaarde.
\end{stappen}}

%========================================================================================
\opdrachtenkop

\begin{oefeningen}
\needspace{10\baselineskip}
\oef{h5-o1}

\needspace{10\baselineskip}
\oef{h5-o2}

\needspace{10\baselineskip}
\oef{h5-o3}

\needspace{10\baselineskip}
\oef{h5-o4}

\needspace{10\baselineskip}
\oef{h5-o5}

\needspace{10\baselineskip}
\oef{h5-o6}

\needspace{10\baselineskip}
\oef{h5-o7}

\needspace{10\baselineskip}
\oef{h5-o8}

\needspace{10\baselineskip}
\oef{h5-o9}

\needspace{10\baselineskip}
\oef{h5-o10}

\needspace{10\baselineskip}
\oef{h5-o11}

\needspace{10\baselineskip}
\oef{h5-o12}
\end{oefeningen}

%----------------------------------------------------------------------------------------
\needspace{14\baselineskip}
\opdrachtsectie{Test jezelf}

\begin{oefeningen}
\oef{h5-t1}

\oef{h5-t2}

\oef{h5-t3}

\oef{h5-t4}

\oef{h5-t5}
\end{oefeningen}

\herhaalschema{%
  \fzv{oef:h5-t1}{oef:h5-o2}%
  \fzv{oef:h5-t2}{oef:h5-o3}%
  \fv{oef:h5-t3}{vb:h5-vgl}{oef:h5-o7}%
  \fv{oef:h5-t4}{vb:h5-vgl}{oef:h5-o8}%
  \fzv{oef:h5-t5}{oef:h5-o9}}

\ifzelfstudie
\section*{Maak je eigen samenvatting}
Sluit dit boek en schrijf in je schrift: hoe je een eerstegraadsvergelijking en -ongelijkheid oplost, de product-nulregel, de discriminant met haar drie gevallen, hoe je een vergelijking met breuken aanpakt, en drie klassieke fouten. Vergelijk daarna met het kader.
\fi
